\chapter{Inference-Time Technique Prompts} \label{ap1:it-prompts}

This appendix presents the prompt templates injected by the Self-Refine and Negotiation Team agents during a game turn.
These prompts are internal to the agent and are never forwarded to the opponent.


\section{Self-Refine Prompts}
\label{ap1:sr-prompts}

The Self-Refine agent \parencite{madaan_self-refine_2023} introduces a feedback-and-refinement loop in every turn.
After generating an initial draft through a \texttt{think()} call, the agent sends two consecutive prompts---a feedback prompt followed by a refine prompt---on a scratch copy of its conversation, then discards the exchange and commits only the final refined response to its persistent history.
By default, the loop runs for two iterations.

\subsection{Feedback Prompt}

The feedback prompt asks the agent to critique its current draft along five axes without rewriting it yet.

\begin{lstlisting}[caption={Self-Refine feedback prompt.},label={lst:sr-feedback}]
    Before finalizing your response, critique it along these axes. Be specific and actionable -- "make it better" is not useful. If a dimension is fine, say OK and move on. Do not rewrite the response yet.

1. Format: Does your response contain all required tags in the correct order, with valid content (e.g., a parseable <proposed trade>)? Note any missing or malformed tag.

2. Payoff alignment: Given your <my goals> (or <my valuation>), does the proposed trade actually advance your payoff? Would a slightly different split give you more utility while still being plausibly acceptable? Quantify if possible.

3. Opponent plausibility: Given what the other player has said or offered so far, is your proposal likely to be accepted, or will it almost certainly be rejected? A rejected proposal wastes your turn.
   
4. Consistency: Is your <reason> consistent with your <player answer> and <newly proposed trade>? Does your <message> contradict your actual move?
   
5. Rule compliance: Are you respecting the proposal budget and turn constraints stated in the rules?

Emit your critique inside <critique> ... </critique> tags. Include at most one concrete, actionable change per axis (e.g. "increase X by 2", "switch from ACCEPT to propose", "remove the contradictory claim in <message>"), not vague suggestions.
\end{lstlisting}

\subsection{Refine Prompt}

After the critique has been generated, the refine prompt instructs the agent to rewrite its full response incorporating the flagged changes.

\begin{lstlisting}[caption={Self-Refine refine prompt.},label={lst:sr-refine}]
    Now rewrite your full response, incorporating the critique. Keep what was fine; change only what the critique flagged. Emit the complete response in the required format (all tags in order, as specified at the start of the game), not a diff. Do not include the critique in the final answer.
\end{lstlisting}


\section{Negotiation Team Prompts}
\label{ap1:team-prompts}

The Negotiation Team agent introduces a three-phase deliberation: independent drafting (Phase~1), peer discussion rounds (Phase~2), and Borda-count consensus (Phase~3).
Phase~1 uses each member's standard game system prompt without modification. The three prompts below drive Phases~2 and~3.

\subsection{Discussion Prompt (Phase~2)}

Each team member receives this prompt once per discussion round. The placeholders \texttt{[own draft]} and \texttt{[peer block]} are substituted at runtime. The peer block lists each distinct peer proposal exactly once, annotated with a support count when multiple members submitted the same text (e.g.\ ``supported by 2 members'').

\begin{lstlisting}[caption={Negotiation Team discussion prompt (Phase~2).},label={lst:team-discussion}]
    Your team is deliberating privately before sending one move. This discussion is never shown to the  other player.

    Your current proposal:
    [own draft]

    Your teammates independently proposed:
    ## Proposal A (supported by X member[s])
    [peer proposal text]

    (additional proposals follow in the same format)

    Weigh these against your own. If a teammate's reasoning is stronger for the team's payoff, adopt or merge it; if yours is stronger, keep it. Then output your (possibly revised) full move in exactly the required response format; all tags in order, as specified at the start of the game. Output only that move.
\end{lstlisting}

\subsection{Ranking Prompt (Phase~3)}

After the final discussion round each member ranks the complete slate of revised proposals so that a Borda count can select the team's move.
The placeholder \texttt{[slate block]} is the set of labelled candidate moves; \texttt{[labels]} is the comma-separated list of candidate letters (e.g.\ A, B, C for a three-member team).

\begin{lstlisting}[caption={Negotiation Team ranking prompt (Phase~3).},label={lst:team-ranking}]
    Your team must pick one of the following candidate moves to send. This vote is private.

    ## Candidate A 
    [candidate A text]

    ## Candidate B 
    [candidate B text]

    (additional candidates follow in the same format)

    Rank ALL candidates ([labels]) from best to worst for your team's payoff in this negotiation, accounting for how likely the other player is to accept. Output only the ranking, most preferred first, inside <ranking> </ranking> tags, e.g. <ranking> A > B > C </ranking>.
\end{lstlisting}

\subsection{Reason-Laundering Prompt (Phase~3)}

Once the Borda count selects the winning candidate, the winning member rewrites only its \texttt{<reason>} block to remove all references to the deliberation (teammates, candidate labels, rankings) that would be invisible on future turns after the conversation is rolled back.
The placeholder \texttt{[original reason]} is the raw \texttt{<reason>} text extracted from the winning draft.

\begin{lstlisting}[caption={Negotiation Team reason-laundering prompt (Phase~3).},label={lst:team-reason}]
Your team has agreed on this move. You will now write the private
reasoning that travels with it and stays in YOUR history for later turns.
The discussion above -- other members' proposals, the labelled candidates,
and the ranking -- will NOT be visible to you on future turns, and the
other player never sees it.

Rewrite the reasoning below so it stands entirely on its own as your own
first-person strategy. Keep the same strategy and the same conclusion. Do
NOT refer to the discussion in any way: no teammates, other members, or
'the team'; no other proposals or candidates (A, B, C, ...); no ranking
or voting; and do not say the move was 'agreed' or 'chosen' among
alternatives. Write in the first person ('I'), as a single negotiator.

Reasoning to rewrite:
[original reason]

Output ONLY the rewritten reasoning, inside <reason> </reason> tags.
\end{lstlisting}
